\documentclass{article}

\usepackage{amssymb,amsmath,tikz,tcolorbox}

\usetikzlibrary{matrix,backgrounds,patterns,positioning}

\setlength{\parindent}{0pt}
\setlength{\parskip}{1em}

\newcommand{\bN}{\mathbb{N}}
\newcommand{\bP}{\mathbb{P}}
\newcommand{\bZ}{\mathbb{Z}}
\newcommand{\bQ}{\mathbb{Q}}
\newcommand{\bR}{\mathbb{R}}
\newcommand{\bC}{\mathbb{C}}
\title{M\"obius Inversion}
\author{Hiten Tandon}
\date{}

\begin{document}
\maketitle

\section{Notation}

\begin{description}
  \item[$\bN$]= Natural Numbers
  \item[$\bP$]= Positive Integers
  \item[$\bZ$]= Integers
  \item[$\bR$]= Real Numbers
  \item[$\bC$]= Complex Numbers
  \item[R]= Commutative Unital RING
  \item[{$[a, b]$}]= $\{x \in \bZ \left| a \le x \le b\right.\}$
  \item[{$[a]$}]= $[1, a]$
\end{description}

\section{M\"obius Function}

Consider two functions $f, g \colon \bN \longrightarrow \bC$ such that
$$g(x) = \sum\limits_{y \le x} f(y)$$

Then, $f$ can be written in terms of $g$ as:
\[
  f(x) =
  \begin{cases}
    g(x) & x = 0 \\
    g(x) - g(x - 1) & x \neq 0
  \end{cases}
\]

This can be proven trivially by expanding the defintion of $g(x)$ in the above relation.

Trivially,
$$g(0) = \sum_{x \le 0} f(x) = f(0)$$
and,
$$g(x) - g(x - 1) = \sum_{y \le x} f(y) - \sum_{z \le x - 1} f(z) = f(x)$$

Crucially, it doesn't matter, how $f$ and $g$ are defined independently of each other.
Just because the relation of $g(x) = \sum\limits_{y \le x} f(y)$ exists,
no matter how else they are originally defined (\S~\ref{sec:def}), we can go always go from computation of one to another
using the same construction.


Also note that, this can be re-written as:
$$f(x) = \sum_{y \le x} \mu(y, x) \cdot g(y)$$
where we define $\mu : \bP \times \bP \longrightarrow \{-1, 0, 1\}$ as:
\[
  \mu(x, y) = \begin{cases}
    -1 & y = x + 1\\
    1 & y = x\\
    0 & otherwise.
  \end{cases}
\]

With this definition of $\mu$ I can now re-write the relation from above as:
$$\boxed{\forall f, g \bP \longrightarrow \bC, f(x) = \sum\limits_{y \le x} g(y). g(x) = \sum\limits_{y \le x} \mu(y, x) \cdot f(y)}$$

Similarly, consider the example where $f, g\colon \bN \times \bN \longrightarrow \bC$ and,
$$g(x, y) = \sum_{z \le x} \sum_{w \le y} f(z, w)$$
The relation can be re-written with $f$ in terms of $g$ as:
\[
  f(x, y) =
  \begin{cases}
    g(x, y) & x = 0, y = 0 \\
    g(x, y) - g(x - 1, y) & x \neq 0, y = 0\\
    g(x, y) - g(x, y - 1) & x = 0, y \neq 0\\
    g(x, y) - g(x - 1, y) - g(x, y - 1) + g(x - 1, y - 1) & x \neq 0, y \neq 0
  \end{cases}
\]

Or, if we take some more liberty we can also do:

\[
  f(x, y) =
  \begin{cases}
    0 & x \notin \bN \text{ or } y \notin \bN \\
    g(x, y) - g(x - 1, y) - g(x, y - 1) + g(x - 1, y - 1) & otherwise.
  \end{cases}
\]

Similar to the previous function, this one can also be re-written as:
$$f(x, y) = \sum_{x' \le x}\sum_{y' \le y} \mu(x, y, x', y') \cdot g(x', y')$$

Where this time we define $\mu : \bP^4 \longrightarrow \{-1, 0, 1\}$ as
\[
  \mu(x, y, x', y') = \begin{cases}
    0 & x' < x - 1 \text{ or } y' < y - 1\\
    -1 & x' = x - 1 \text { xor } y' = y - 1\\
    1 & otherwise.
  \end{cases} 
\]

You can take your time to verify the validity of this claim.\\
Once you do, notice that just like the last time, this time as well\\
all we really need is the relation between $f$ and $g$. We don't particularly\\
care about the definitions following any rules. Just this relation holding is\\
enough to jump from one to other.

\begin{tcolorbox}[colframe=red!50, title=Fun tidbit]
If we call $\mu$ we defined in the first example $\mu_1$ and the one from the second example $\mu_2$
Then observe that we can rewrite $\mu_2$ in terms of $\mu_1$ as

$$\mu_2(x, y, x', y') = \mu_1(x, x') \cdot \mu_1(y, y')$$
\end{tcolorbox}

Now for the most important example,
$\forall f, g: \bP \longrightarrow \bC$
$$f(x) = \sum_{d | x} g(d)$$
can be re-written as
$$g(x) = \sum_{d | x} \mu(x, d) \cdot f(d)$$
where $\mu : \bP^2 \longrightarrow \{-1, 0, 1\}$
is defined as
\[
 \mu(x, y) =
  \begin{cases}
    (-1)^{n\left(p \in \text{Primes} \colon  p | \frac x y \right)} & \forall p \in \bP, p^2 \nmid \frac x y.\\
    0 & otherwise
  \end{cases} 
\]

\section{Formal definition and constraints of a M\"obius Function}\label{sec:def}
\subsection{Partially Ordered Sets / Posets}

A partially ordered set also known as a Poset is a pair of a set $P$ and a relation $\le \colon P^2$ s.t it staisfies:
\begin{enumerate}
\item $\forall x \in P. x \le x$ (Reflexivity)
\item $\forall x, y, z \in P, x \le y, y \le z. x \le z.$ (Transitivity)
\item $\forall x, y \in P, x \le y, y \le x. x = y.$ (Anti-symmetry)
\end{enumerate}

Common examples of Posets include
\begin{enumerate}
\item all subsets of $(\bC, \le)$
\item $(\bP, \mid)$
\item $\forall a, b \in \bZ. [a,b]$
\item $\forall X: \text{Set}. (2^X, \subseteq)$
\item $\left(\bR^2, \left\{(x, y), (x', y') \in \bR^2 \left| x \le x', y \le y'\right.\right\}\right)$
\end{enumerate}

\begin{tcolorbox}[colframe=blue!50, title=Chains and Anti-chains]
A Poset $(P, \le)$ is called a chain if $\forall x, y \in P, x \le y \text{ or } y \le x$.
Common examples of this would include $\bN, \bP, \bR$ and their subsets.
\tcblower
A Poset $(P, \le)$ is called an anti-chain if $\forall x, y \in P, x \le y \iff y \le x$.
A common example of an antichain is \{German, French, Italian\}
\end{tcolorbox}

\subsection{Hasse Diagrams}
A common way to represent a Poset is by using Hasse diagrams.
A Hasse diagram is a graphical representation of a Poset where the following conditions are met:
\begin{enumerate}
\item Vertices of the graph represent an element of the poset $P$
\item There is an edge from $x$ to $y$ if $x$ covers $y$.
\end{enumerate}

Cover here means the following:
$\forall x, y \in P. \nexists z \in P, z \leq x, y \leq z, x \neq z, y \neq z.$
Then $x$ is said to cover $y$.

\end{document}
